\chapter{THIẾT KẾ CƠ SỞ DỮ LIỆU}

\section{Thực thể và mối quan hệ}

Từ phân tích ở Chương 1, cơ sở dữ liệu được chia thành sáu nhóm với 30 bảng
(Bảng~\ref{tab:groups}). Ngoài tám bảng nhật ký và đánh giá, có 22 bảng chính, vượt yêu cầu
tối thiểu 6 bảng của đề bài.

\begin{table}[H]
\centering
\caption{Các nhóm bảng trong cơ sở dữ liệu \texttt{video\_archive}}
\label{tab:groups}
\renewcommand{\arraystretch}{1.2}
\begin{tabularx}{\textwidth}{>{\raggedright\arraybackslash}p{3.8cm} >{\raggedright\arraybackslash}X}
\toprule
\textbf{Nhóm} & \textbf{Bảng} \\
\midrule
Danh mục thư viện & \texttt{organizations}, \texttt{collections}, \texttt{topics}, \texttt{video\_topics} \\
Kho dữ liệu video & \texttt{videos}, \texttt{keyframes}, \texttt{ocr\_texts}, \texttt{asr\_segments} \\
Chỉ mục văn bản & \texttt{stopwords}, \texttt{terms}, \texttt{ocr\_postings}, \texttt{asr\_postings} \\
Tài khoản và phân quyền & \texttt{users}, \texttt{roles}, \texttt{permissions}, \texttt{user\_roles}, \texttt{role\_permissions}, \texttt{sessions} \\
Yêu cầu khai thác và cấp phép & \texttt{requests}, \texttt{ground\_truth}, \texttt{selections}, \texttt{usage\_licenses} \\
Nhật ký và đánh giá & \texttt{sources}, \texttt{audit\_log}, \texttt{search\_log}, \texttt{search\_log\_sources}, \texttt{eval\_runs}, \texttt{eval\_run\_sources}, \texttt{eval\_requests}, \texttt{eval\_results} \\
\bottomrule
\end{tabularx}
\end{table}

Các mối quan hệ chính và bản số của chúng được liệt kê ở Bảng~\ref{tab:relations}.
Quan hệ nhiều--nhiều được tách thành bảng trung gian có khoá chính ghép.

{\renewcommand{\arraystretch}{1.15}
\begin{longtable}{>{\raggedright\arraybackslash}p{5.4cm} >{\centering\arraybackslash}p{1.4cm} >{\raggedright\arraybackslash}p{7.0cm}}
\caption{Các mối quan hệ chính}\label{tab:relations}\\
\toprule \textbf{Quan hệ} & \textbf{Bản số} & \textbf{Ý nghĩa} \\ \midrule \endfirsthead
\toprule \textbf{Quan hệ} & \textbf{Bản số} & \textbf{Ý nghĩa} \\ \midrule \endhead
\bottomrule \endfoot
organizations -- collections & 1 : N & Một đơn vị đóng góp nhiều bộ sưu tập \\
collections -- videos & 1 : N & Một bộ sưu tập có nhiều video; video thuộc đúng một bộ sưu tập \\
topics -- topics & 1 : N & Quan hệ đệ quy: chủ đề cha có nhiều chủ đề con \\
videos -- topics (qua video\_topics) & N : N & Một video có nhiều chủ đề, một chủ đề gắn cho nhiều video \\
videos -- keyframes & 1 : N & Một video có nhiều khung hình \\
keyframes -- ocr\_texts & 1 : 0..1 & Một khung hình có thể có chữ \\
videos -- asr\_segments & 1 : N & Một video có nhiều đoạn lời nói \\
terms -- keyframes (qua ocr\_postings) & N : N & Từ khoá xuất hiện ở nhiều khung hình \\
terms -- asr\_segments (qua asr\_postings) & N : N & Từ khoá xuất hiện ở nhiều đoạn lời nói \\
organizations -- users & 1 : N & Bạn đọc, cán bộ thuộc một đơn vị (hoặc là cá nhân) \\
users -- requests & 1 : N & Một bạn đọc gửi nhiều yêu cầu khai thác \\
requests -- keyframes (qua selections) & N : N & Một yêu cầu có nhiều đoạn tư liệu, một khung hình dùng cho nhiều yêu cầu \\
requests -- usage\_licenses & 1 : 0..1 & Yêu cầu được duyệt có đúng một giấy phép \\
organizations -- usage\_licenses & 1 : N & Đơn vị được cấp phép \\
users -- usage\_licenses & 1 : N & Cán bộ cấp phép \\
requests -- ground\_truth & 1 : 0..1 & Chỉ yêu cầu mẫu mới có vị trí đã xác minh \\
users -- roles (qua user\_roles) & N : N & Một người có nhiều vai trò \\
roles -- permissions (qua role\_permissions) & N : N & Một vai trò có nhiều quyền \\
users -- sessions & 1 : N & Một người có nhiều phiên đăng nhập \\
search\_log / eval\_runs -- sources & N : N & Một lượt tìm/đánh giá dùng nhiều nguồn \\
eval\_runs -- eval\_requests -- eval\_results & 1 : N : N & Lượt đánh giá $\rightarrow$ yêu cầu $\rightarrow$ kết quả xếp hạng \\
\end{longtable}}

\section{Sơ đồ thực thể -- liên kết}

Do có 30 bảng, sơ đồ ERD được tách theo nhóm (Hình~\ref{fig:erd_catalog}--\ref{fig:erd_logs}).
Thuộc tính gạch chân là khoá chính, \textit{(FK)} là khoá ngoại, \textit{(UQ)} là duy nhất;
khung nét đứt là bảng thuộc nhóm khác, chỉ vẽ lại để thấy liên kết. Nhãn trên đường nối là
bản số.

\begin{figure}[H]
\centering
\resizebox{\textwidth}{!}{%
\begin{tikzpicture}
\node[entity] (org) at (0,0) {\textbf{organizations}\nodepart{second}\attrs{\pk{organization\_id} \\ org\_code (UQ) \\ org\_name (UQ) \\ org\_type \\ created\_at}};
\node[entity] (col) at (5,0) {\textbf{collections}\nodepart{second}\attrs{\pk{collection\_id} \\ collection\_code (UQ) \\ title \\ organization\_id (FK) \\ access\_level \\ description, created\_at}};
\node[entity] (vid) at (10,0) {\textbf{videos}\nodepart{second}\attrs{\pk{video\_id} \\ batch (tính toán) \\ collection\_id (FK) \\ title \\ fps}};
\node[entity] (vt) at (10,-4.6) {\textbf{video\_topics}\nodepart{second}\attrs{\pk{video\_id} (FK) \\ \pk{topic\_id} (FK)}};
\node[entity] (top) at (5,-4.6) {\textbf{topics}\nodepart{second}\attrs{\pk{topic\_id} \\ topic\_name (UQ) \\ parent\_topic\_id (FK)}};
\node[refentity] (kf) at (14.6,0) {\textbf{keyframes}\nodepart{second}\attrs{\pk{keyframe\_id} \\ video\_id (FK) \\ \dots}};
\draw[rel] (org) -- node[card, pos=0.2]{1} node[card, pos=0.8]{N} (col);
\draw[rel] (col) -- node[card, pos=0.2]{1} node[card, pos=0.8]{N} (vid);
\draw[rel] (vid) -- node[card, pos=0.2]{1} node[card, pos=0.8]{N} (vt);
\draw[rel] (top) -- node[card, pos=0.2]{1} node[card, pos=0.8]{N} (vt);
\draw[rel] (vid) -- node[card, pos=0.2]{1} node[card, pos=0.8]{N} (kf);
\draw[rel] (top.west) to[out=180, in=-90, looseness=1.4] node[card, pos=0.15]{N}
    ($(top.west)+(-1.6,0)$) to[out=90, in=150, looseness=1.4] node[card, pos=0.85]{1} (top.north west);
\node[font=\scriptsize\itshape, align=center] at ($(top.west)+(-2.4,0)$) {cha --\\con};
\end{tikzpicture}}
\caption{ERD nhóm danh mục thư viện}
\label{fig:erd_catalog}
\end{figure}

\begin{figure}[H]
\centering
\resizebox{\textwidth}{!}{%
\begin{tikzpicture}
\node[refentity] (videos) at (0,0) {\textbf{videos}\nodepart{second}\attrs{\pk{video\_id} \\ \dots}};
\node[entity] (keyframes) at (4.4,0) {\textbf{keyframes}\nodepart{second}\attrs{\pk{keyframe\_id} \\ video\_id (FK) \\ sample\_n \\ pts\_time \\ frame\_idx}};
\node[entity] (ocr) at (8.8,0) {\textbf{ocr\_texts}\nodepart{second}\attrs{\pk{keyframe\_id} (FK) \\ ocr\_text \\ text\_folded (tính toán)}};
\node[entity] (ocrp) at (13.2,0) {\textbf{ocr\_postings}\nodepart{second}\attrs{\pk{term\_id} (FK) \\ \pk{keyframe\_id} (FK)}};
\node[entity] (asr) at (0,-4.2) {\textbf{asr\_segments}\nodepart{second}\attrs{\pk{segment\_id} \\ video\_id (FK) \\ start\_sec, end\_sec \\ asr\_text \\ text\_folded (tính toán)}};
\node[entity] (asrp) at (4.4,-4.2) {\textbf{asr\_postings}\nodepart{second}\attrs{\pk{term\_id} (FK) \\ \pk{segment\_id} (FK)}};
\node[entity] (terms) at (8.8,-4.2) {\textbf{terms}\nodepart{second}\attrs{\pk{term\_id} \\ term (UQ)}};
\node[entity] (stop) at (13.2,-4.2) {\textbf{stopwords}\nodepart{second}\attrs{\pk{word}}};
\draw[rel] (videos) -- node[card, pos=0.18]{1} node[card, pos=0.82]{N} (keyframes);
\draw[rel] (keyframes) -- node[card, pos=0.2]{1} node[card, pos=0.8]{0..1} (ocr);
\draw[rel] (ocr) -- node[card, pos=0.2]{1} node[card, pos=0.8]{N} (ocrp);
\draw[rel] (videos) -- node[card, pos=0.2]{1} node[card, pos=0.8]{N} (asr);
\draw[rel] (asr) -- node[card, pos=0.2]{1} node[card, pos=0.8]{N} (asrp);
\draw[rel] (terms) -- node[card, pos=0.2]{1} node[card, pos=0.8]{N} (asrp);
\draw[rel] (terms) -- node[card, pos=0.2]{1} node[card, pos=0.8]{N} (ocrp);
\end{tikzpicture}}
\caption{ERD nhóm kho dữ liệu video và chỉ mục văn bản}
\label{fig:erd_corpus}
\end{figure}

\begin{figure}[H]
\centering
\resizebox{\textwidth}{!}{%
\begin{tikzpicture}
\node[entity] (gt) at (-5.4,0) {\textbf{ground\_truth}\nodepart{second}\attrs{\pk{request\_id} (FK) \\ video\_id (FK) \\ start\_sec, end\_sec}};
\node[entity] (q) at (0,0) {\textbf{requests}\nodepart{second}\attrs{\pk{request\_id} \\ request\_text \\ origin \\ requested\_by (FK) \\ purpose \\ status \\ created\_at}};
\node[entity] (a) at (5.4,0) {\textbf{selections}\nodepart{second}\attrs{\pk{selection\_id} \\ request\_id (FK) \\ keyframe\_id (FK) \\ user\_id (FK) \\ note, created\_at}};
\node[refentity] (k) at (10.4,0) {\textbf{keyframes}\nodepart{second}\attrs{\pk{keyframe\_id} \\ \dots}};
\node[refentity] (v) at (-5.4,-6) {\textbf{videos}\nodepart{second}\attrs{\pk{video\_id} \\ \dots}};
\node[entity] (l) at (0,-6) {\textbf{usage\_licenses}\nodepart{second}\attrs{\pk{license\_id} \\ request\_id (FK, UQ) \\ licensee\_org\_id (FK) \\ granted\_by (FK) \\ granted\_at \\ valid\_from, valid\_to \\ conditions}};
\node[refentity] (u) at (5.4,-6) {\textbf{users}\nodepart{second}\attrs{\pk{user\_id} \\ \dots}};
\node[refentity] (o) at (0,-11.4) {\textbf{organizations}\nodepart{second}\attrs{\pk{organization\_id} \\ \dots}};
\draw[rel] (q) -- node[card, pos=0.2]{1} node[card, pos=0.8]{N} (a);
\draw[rel] (k) -- node[card, pos=0.2]{1} node[card, pos=0.8]{N} (a);
\draw[rel] (u) -- node[card, pos=0.2]{1} node[card, pos=0.8]{N} (a);
\draw[rel] (q) -- node[card, pos=0.2]{1} node[card, pos=0.8]{0..1} (gt);
\draw[rel] (v) -- node[card, pos=0.2]{1} node[card, pos=0.8]{N} (gt);
\draw[rel] (q) -- node[card, pos=0.2]{1} node[card, pos=0.8]{0..1} (l);
\draw[rel] (u) -- node[card, pos=0.2]{1} node[card, pos=0.8]{N} node[font=\scriptsize\itshape, above, pos=0.5]{cấp phép} (l);
\draw[rel] (u) -- node[card, pos=0.2]{1} node[card, pos=0.8]{N} node[font=\scriptsize\itshape, sloped, above, pos=0.45]{gửi yêu cầu} (q);
\draw[rel] (o) -- node[card, pos=0.2]{1} node[card, pos=0.8]{N} (l);
\end{tikzpicture}}
\caption{ERD nhóm yêu cầu khai thác và cấp phép}
\label{fig:erd_requests}
\end{figure}

\begin{figure}[H]
\centering
\resizebox{0.95\textwidth}{!}{%
\begin{tikzpicture}
\node[entity] (u) at (0,0) {\textbf{users}\nodepart{second}\attrs{\pk{user\_id} \\ username (UQ) \\ full\_name \\ password\_hash \\ organization\_id (FK) \\ user\_type \\ is\_active \\ created\_at, last\_login\_at}};
\node[entity] (ur) at (4.5,0) {\textbf{user\_roles}\nodepart{second}\attrs{\pk{user\_id} (FK) \\ \pk{role\_id} (FK)}};
\node[entity] (r) at (8.6,0) {\textbf{roles}\nodepart{second}\attrs{\pk{role\_id} \\ role\_name (UQ) \\ description}};
\node[entity] (rp) at (12.8,0) {\textbf{role\_permissions}\nodepart{second}\attrs{\pk{role\_id} (FK) \\ \pk{permission\_code} (FK)}};
\node[entity] (p) at (12.8,-4.2) {\textbf{permissions}\nodepart{second}\attrs{\pk{permission\_code} \\ description}};
\node[entity] (s) at (4.5,-4.2) {\textbf{sessions}\nodepart{second}\attrs{\pk{session\_hash} \\ user\_id (FK) \\ created\_at, expires\_at}};
\node[refentity] (o) at (0,-4.6) {\textbf{organizations}\nodepart{second}\attrs{\pk{organization\_id} \\ \dots}};
\draw[rel] (u) -- node[card, pos=0.2]{1} node[card, pos=0.8]{N} (ur);
\draw[rel] (r) -- node[card, pos=0.2]{1} node[card, pos=0.8]{N} (ur);
\draw[rel] (r) -- node[card, pos=0.2]{1} node[card, pos=0.8]{N} (rp);
\draw[rel] (p) -- node[card, pos=0.2]{1} node[card, pos=0.8]{N} (rp);
\draw[rel] (u) -- node[card, pos=0.2]{1} node[card, pos=0.8]{N} (s);
\draw[rel] (o) -- node[card, pos=0.25]{0..1} node[card, pos=0.8]{N} (u);
\end{tikzpicture}}
\caption{ERD nhóm tài khoản và phân quyền}
\label{fig:erd_rbac}
\end{figure}

\begin{figure}[H]
\centering
\resizebox{0.95\textwidth}{!}{%
\begin{tikzpicture}
\node[refentity] (u) at (0,0) {\textbf{users}\nodepart{second}\attrs{\pk{user\_id} \\ \dots}};
\node[entity] (al) at (4.5,1.6) {\textbf{audit\_log}\nodepart{second}\attrs{\pk{log\_id} \\ user\_id (FK) \\ action, detail \\ created\_at}};
\node[entity] (sl) at (4.5,-1.6) {\textbf{search\_log}\nodepart{second}\attrs{\pk{search\_id} \\ user\_id (FK) \\ request\_text, top\_k \\ result\_count, latency\_ms}};
\node[entity] (sls) at (9,-1.6) {\textbf{search\_log\_sources}\nodepart{second}\attrs{\pk{search\_id} (FK) \\ \pk{source\_code} (FK)}};
\node[entity] (src) at (13.4,-1.6) {\textbf{sources}\nodepart{second}\attrs{\pk{source\_code} \\ description}};
\node[entity] (ers) at (13.4,-5.4) {\textbf{eval\_run\_sources}\nodepart{second}\attrs{\pk{run\_id} (FK) \\ \pk{source\_code} (FK)}};
\node[entity] (er) at (9,-5.4) {\textbf{eval\_runs}\nodepart{second}\attrs{\pk{run\_id} \\ run\_name, top\_k \\ tolerance\_sec}};
\node[entity] (eq) at (4.5,-5.4) {\textbf{eval\_requests}\nodepart{second}\attrs{\pk{run\_id} (FK) \\ \pk{request\_id} (FK) \\ latency\_ms}};
\node[refentity] (g) at (0,-5.4) {\textbf{ground\_truth}\nodepart{second}\attrs{\pk{request\_id} \\ \dots}};
\node[entity] (res) at (4.5,-9.2) {\textbf{eval\_results}\nodepart{second}\attrs{\pk{run\_id}, \pk{request\_id} (FK) \\ \pk{rank\_no} \\ keyframe\_id (FK)}};
\node[refentity] (k) at (0,-9.2) {\textbf{keyframes}\nodepart{second}\attrs{\pk{keyframe\_id} \\ \dots}};
\draw[rel] (u) -- node[card, pos=0.2]{1} node[card, pos=0.8]{N} (al);
\draw[rel] (u) -- node[card, pos=0.2]{1} node[card, pos=0.8]{N} (sl);
\draw[rel] (sl) -- node[card, pos=0.2]{1} node[card, pos=0.8]{N} (sls);
\draw[rel] (src) -- node[card, pos=0.2]{1} node[card, pos=0.8]{N} (sls);
\draw[rel] (src) -- node[card, pos=0.2]{1} node[card, pos=0.8]{N} (ers);
\draw[rel] (er) -- node[card, pos=0.2]{1} node[card, pos=0.8]{N} (ers);
\draw[rel] (er) -- node[card, pos=0.2]{1} node[card, pos=0.8]{N} (eq);
\draw[rel] (g) -- node[card, pos=0.2]{1} node[card, pos=0.8]{N} (eq);
\draw[rel] (eq) -- node[card, pos=0.2]{1} node[card, pos=0.8]{N} (res);
\draw[rel] (k) -- node[card, pos=0.2]{1} node[card, pos=0.8]{N} (res);
\end{tikzpicture}}
\caption{ERD nhóm nhật ký và đánh giá}
\label{fig:erd_logs}
\end{figure}

\section{Chuẩn hoá đến dạng chuẩn 3}

\subsection{Phụ thuộc hàm của các bảng chính}

Lược đồ được kiểm tra theo định nghĩa của Codd~\cite{codd1970relational} và các dạng chuẩn
trong~\cite{elmasri2016fundamentals}: một lược đồ đạt 3NF nếu với mọi phụ thuộc hàm không
tầm thường $X \rightarrow A$, hoặc $X$ là siêu khoá, hoặc $A$ thuộc một khoá ứng viên.
Các phụ thuộc hàm chính:
\begin{itemize}
    \item \texttt{organizations}: \texttt{organization\_id} $\rightarrow$ mọi thuộc tính;
    \texttt{org\_code} và \texttt{org\_name} cũng là khoá ứng viên.
    \item \texttt{collections}: \texttt{collection\_id} $\rightarrow$ (\texttt{collection\_code},
    \texttt{title}, \texttt{organization\_id}, \texttt{access\_level}, \dots);
    \texttt{collection\_code} là khoá ứng viên.
    \item \texttt{topics}: \texttt{topic\_id} $\rightarrow$ (\texttt{topic\_name},
    \texttt{parent\_topic\_id}); \texttt{topic\_name} là khoá ứng viên.
    \item \texttt{videos}: \texttt{video\_id} $\rightarrow$ (\texttt{collection\_id},
    \texttt{title}, \texttt{fps}, \texttt{batch}).
    \item \texttt{keyframes}: có bốn khoá ứng viên: \texttt{keyframe\_id},
    (\texttt{video\_id}, \texttt{sample\_n}), (\texttt{video\_id}, \texttt{pts\_time}) và
    (\texttt{video\_id}, \texttt{frame\_idx}). Mỗi khoá xác định mọi thuộc tính còn lại; dữ
    liệu thật không có dòng trùng ở bất kỳ khoá nào, và cả bốn đều được khai báo bằng
    ràng buộc PRIMARY KEY hoặc UNIQUE.
    \item \texttt{ocr\_texts}: \texttt{keyframe\_id} $\rightarrow$ \texttt{ocr\_text};
    \texttt{ocr\_text} $\rightarrow$ \texttt{text\_folded} (bỏ dấu là hàm của văn bản).
    Tương tự với \texttt{asr\_segments}.
    \item \texttt{users}: \texttt{user\_id} $\rightarrow$ mọi thuộc tính;
    \texttt{username} là khoá ứng viên.
    \item \texttt{requests}: \texttt{request\_id} $\rightarrow$ (\texttt{request\_text},
    \texttt{origin}, \texttt{requested\_by}, \texttt{purpose}, \texttt{status},
    \texttt{created\_at}).
    \item \texttt{selections}: \texttt{selection\_id} $\rightarrow$ mọi thuộc tính;
    (\texttt{request\_id}, \texttt{keyframe\_id}) là khoá ứng viên.
    \item \texttt{usage\_licenses}: \texttt{license\_id} $\rightarrow$ mọi thuộc tính;
    \texttt{request\_id} là khoá ứng viên (mỗi yêu cầu tối đa một giấy phép).
    \item \texttt{eval\_results}: (\texttt{run\_id}, \texttt{request\_id}, \texttt{rank\_no})
    $\rightarrow$ \texttt{keyframe\_id}; (\texttt{run\_id}, \texttt{request\_id},
    \texttt{keyframe\_id}) cũng là khoá ứng viên.
\end{itemize}

\subsection{Các bước chuẩn hoá}

Bản thiết kế đầu tiên của phần kho video chưa đạt 3NF; nhóm kiểm tra phụ thuộc hàm trên dữ
liệu thật và sửa theo từng bước. Khi mở rộng sang danh mục thư viện và quy trình cấp phép,
nhóm cũng cân nhắc một số phương án đơn giản hơn nhưng vi phạm chuẩn, rồi loại bỏ chúng
(Bảng~\ref{tab:normalization}).

{\renewcommand{\arraystretch}{1.25}
\begin{longtable}{>{\raggedright\arraybackslash}p{1.4cm} >{\raggedright\arraybackslash}p{6.5cm} >{\raggedright\arraybackslash}p{6.5cm}}
\caption{Các thay đổi trong quá trình chuẩn hoá}\label{tab:normalization}\\
\toprule
\textbf{Bước} & \textbf{Vấn đề} & \textbf{Cách sửa} \\
\midrule \endfirsthead
\toprule
\textbf{Bước} & \textbf{Vấn đề} & \textbf{Cách sửa} \\
\midrule \endhead
\bottomrule \endfoot
1NF & \texttt{search\_log.sources} và \texttt{eval\_runs.sources} lưu danh sách
\texttt{"visual,ocr,asr"} trong một ô, tức thuộc tính đa trị. Không thể ràng buộc giá trị
hợp lệ hay truy vấn ``những lượt dùng OCR'' bằng phép so sánh bằng. &
Thêm bảng danh mục \texttt{sources} và hai bảng trung gian \texttt{search\_log\_sources},
\texttt{eval\_run\_sources} với khoá chính ghép và khoá ngoại tới \texttt{sources}. \\
\addlinespace
1NF & Phương án lưu chủ đề của video thành chuỗi \texttt{"Nước rút, Về đích"} trong
\texttt{videos} lặp lại lỗi trên, và không biểu diễn được cây chủ đề. &
Bảng \texttt{topics} (có khoá ngoại tự tham chiếu tới chủ đề cha) và bảng trung gian
\texttt{video\_topics}. Chuỗi chủ đề chỉ được ghép khi hiển thị, trong view
\texttt{v\_video\_catalog}. \\
\addlinespace
2NF & Kiểm tra các bảng có khoá ghép: \texttt{user\_roles}, \texttt{role\_permissions},
\texttt{ocr\_postings}, \texttt{asr\_postings}, \texttt{video\_topics}, hai bảng nguồn (toàn
khoá); \texttt{eval\_requests}: \texttt{latency\_ms} phụ thuộc cả (\texttt{run\_id},
\texttt{request\_id}); \texttt{eval\_results}: \texttt{keyframe\_id} phụ thuộc cả ba cột khoá. &
Không có phụ thuộc bộ phận, không cần tách. \\
\addlinespace
3NF & \texttt{text\_folded} phụ thuộc bắc cầu qua \texttt{ocr\_text} /
\texttt{asr\_text}, vốn không phải khoá. Khi sửa văn bản mà quên sửa
\texttt{text\_folded} thì dữ liệu mâu thuẫn (dị thường cập nhật). &
Chuyển thành cột tính toán \texttt{PERSISTED} \texttt{AS dbo.fn\_fold\_vi(ocr\_text)}.
Giá trị do hệ quản trị tự tính và lưu, không ai ghi trực tiếp được, nên phụ thuộc
được chính hệ quản trị bảo đảm. Nhóm giữ giá trị lưu thay vì bỏ cột vì cần dùng nó khi
lập chỉ mục. \\
\addlinespace
3NF & Phương án lưu đơn vị sở hữu ngay trên video (\texttt{videos.organization\_id}) cho
tiện truy vấn. Khi đó \texttt{video\_id} $\rightarrow$ \texttt{collection\_id}
$\rightarrow$ \texttt{organization\_id} là phụ thuộc bắc cầu; chuyển bộ sưu tập sang đơn
vị khác sẽ phải sửa mọi video của nó. &
Không lưu; đơn vị sở hữu suy ra qua \texttt{collections}. View \texttt{v\_video\_catalog}
ghép sẵn cho giao diện. \\
\addlinespace
3NF & Bản đầu định lưu sẵn khung hình gần nhất của mỗi đoạn lời nói
(\texttt{asr\_segments.keyframe\_id}). Giá trị này suy ra từ (\texttt{video\_id},
\texttt{start\_sec}, \texttt{end\_sec}) và dữ liệu của bảng \texttt{keyframes}. &
Không lưu; hàm \texttt{fn\_nearest\_keyframe} tính khi truy vấn bằng hai lần dò chỉ mục. \\
\addlinespace
Dư thừa & Phương án thêm trạng thái ``đang xử lý'' cho yêu cầu đã có đoạn tư liệu được
chọn. Trạng thái này suy ra hoàn toàn từ việc có dòng trong \texttt{selections} hay không. &
Không lưu; \texttt{v\_request\_progress} đếm số đoạn đã chọn. \texttt{status} chỉ giữ ba
giá trị không suy ra được: chờ xử lý, đã cấp phép, từ chối. \\
\addlinespace
Dư thừa & \texttt{keyframes.image\_path} suy ra hoàn toàn từ (\texttt{video\_id},
\texttt{sample\_n}) theo quy ước tên tệp (đúng với 100\% số dòng). Không vi phạm 3NF vì vế
trái là khoá ứng viên, nhưng là dữ liệu thừa. & Bỏ cột; đường dẫn do ứng dụng dựng từ hai
cột khoá. \\
\addlinespace
Dư thừa & \texttt{videos.batch} là tiền tố của \texttt{video\_id}. & Cột tính toán
\texttt{PERSISTED} có chỉ mục. \\
\end{longtable}}

\subsection{Kết quả và dư thừa có kiểm soát}

Sau các bước trên, mọi phụ thuộc hàm không tầm thường trong từng bảng đều có vế trái là
siêu khoá, nên lược đồ đạt 3NF và cả BCNF. Nhóm giữ có chủ đích hai dư thừa \textit{giữa các
bảng} (không vi phạm dạng chuẩn của từng bảng) và kiểm tra chúng bằng script:
\begin{itemize}
    \item \texttt{keyframes.pts\_time} bằng \texttt{frame\_idx / videos.fps}. Trong bảng
    \texttt{keyframes}, cả (\texttt{video\_id}, \texttt{pts\_time}) và (\texttt{video\_id},
    \texttt{frame\_idx}) đều là khoá ứng viên. Nhóm giữ \texttt{pts\_time} vì đây là số đo gốc
    trong dữ liệu và là cột của chỉ mục thời gian \texttt{UQ\_keyframes\_video\_time}. Mục E1
    của \texttt{11\_tests.sql}: 0 dòng lệch trên 5.913 khung hình.
    \item \texttt{requests.status = 'approved'} khi và chỉ khi yêu cầu có dòng trong
    \texttt{usage\_licenses}. Nhóm giữ cột \texttt{status} vì hai giá trị còn lại (chờ xử lý,
    từ chối) không suy ra được từ bảng nào, và vì hàng đợi ``yêu cầu chờ xử lý'' cần chỉ mục
    \texttt{IX\_requests\_status}. Tính nhất quán được bảo đảm bằng thiết kế: ứng dụng không có
    quyền \texttt{UPDATE} trên \texttt{requests}, trạng thái chỉ đổi qua
    \texttt{sp\_grant\_license} (ghi giấy phép và trạng thái trong cùng một giao dịch) hoặc
    \texttt{sp\_reject\_request}. Mục E2: 0 dòng lệch trên 120 yêu cầu.
\end{itemize}

Ngược lại, \texttt{usage\_licenses.licensee\_org\_id} trông giống dữ liệu suy ra được (từ
người gửi yêu cầu sang đơn vị của người đó) nhưng không phải: nó là \textit{ảnh chụp} đơn vị
tại thời điểm cấp phép. Nếu bạn đọc chuyển sang đơn vị khác, giấy phép đã cấp vẫn thuộc đơn
vị cũ; suy ra lại từ \texttt{users.organization\_id} hiện tại sẽ cho kết quả sai. Tương tự,
\texttt{users.user\_type} không suy ra được từ loại đơn vị: một trường đại học có cả nhà
nghiên cứu lẫn sinh viên.

\section{Danh sách bảng và mô tả cột}

Các bảng dưới đây được sinh tự động từ catalog của SQL Server (\texttt{sys.columns},
\texttt{sys.key\_constraints}, \texttt{sys.foreign\_keys}, \texttt{sys.default\_constraints})
bằng \texttt{report/tools/gen\_schema.py}, nên luôn khớp với script \texttt{02\_tables.sql}.
Cột ``Khoá / mặc định'' ghi PK, UQ, FK (kèm bảng được tham chiếu), IDENTITY và giá trị
DEFAULT.

{\small
\subsection{Nhóm danh mục thư viện}

\textbf{Bảng \texttt{organizations}} -- Đơn vị thành viên: đóng góp bộ sưu tập và/hoặc có bạn đọc khai thác.

\setlength{\tabcolsep}{4pt}
\begin{longtable}{>{\raggedright\arraybackslash}p{0.21\textwidth}>{\raggedright\arraybackslash}p{0.19\textwidth}>{\centering\arraybackslash}p{0.07\textwidth}>{\raggedright\arraybackslash}p{0.17\textwidth}>{\raggedright\arraybackslash}p{0.26\textwidth}}
\caption{Cấu trúc bảng \texttt{organizations}}\label{tab:col_organizations}\\
\toprule \textbf{Cột} & \textbf{Kiểu} & \textbf{Null} & \textbf{Khoá / mặc định} & \textbf{Mô tả} \\ \midrule \endfirsthead
\toprule \textbf{Cột} & \textbf{Kiểu} & \textbf{Null} & \textbf{Khoá / mặc định} & \textbf{Mô tả} \\ \midrule \endhead
\bottomrule \endfoot
\texttt{organization\_id} & {\footnotesize\texttt{smallint}} & không & PK, IDENTITY & Mã đơn vị (tự tăng) \\
\texttt{org\_code} & {\footnotesize\texttt{varchar(16)}} & không & UQ & Mã viết tắt, vd. LIB \\
\texttt{org\_name} & {\footnotesize\texttt{nvarchar(200)}} & không & UQ & Tên đơn vị \\
\texttt{org\_type} & {\footnotesize\texttt{varchar(16)}} & không &  & Loại: library, tv, press, school, museum \\
\texttt{created\_at} & {\footnotesize\texttt{datetime2(0)}} & không & DF sysdatetime() & Thời điểm tham gia \\
\end{longtable}

\textbf{Bảng \texttt{collections}} -- Bộ sưu tập do một đơn vị đóng góp và giữ bản quyền.

\setlength{\tabcolsep}{4pt}
\begin{longtable}{>{\raggedright\arraybackslash}p{0.21\textwidth}>{\raggedright\arraybackslash}p{0.19\textwidth}>{\centering\arraybackslash}p{0.07\textwidth}>{\raggedright\arraybackslash}p{0.17\textwidth}>{\raggedright\arraybackslash}p{0.26\textwidth}}
\caption{Cấu trúc bảng \texttt{collections}}\label{tab:col_collections}\\
\toprule \textbf{Cột} & \textbf{Kiểu} & \textbf{Null} & \textbf{Khoá / mặc định} & \textbf{Mô tả} \\ \midrule \endfirsthead
\toprule \textbf{Cột} & \textbf{Kiểu} & \textbf{Null} & \textbf{Khoá / mặc định} & \textbf{Mô tả} \\ \midrule \endhead
\bottomrule \endfoot
\texttt{collection\_id} & {\footnotesize\texttt{int}} & không & PK, IDENTITY & Mã bộ sưu tập (tự tăng) \\
\texttt{collection\_code} & {\footnotesize\texttt{varchar(16)}} & không & UQ & Mã viết tắt, vd. DXD \\
\texttt{title} & {\footnotesize\texttt{nvarchar(200)}} & không &  & Tên bộ sưu tập \\
\texttt{organization\_id} & {\footnotesize\texttt{smallint}} & không & FK $\rightarrow$ organizations & Đơn vị đóng góp, giữ bản quyền \\
\texttt{access\_level} & {\footnotesize\texttt{varchar(16)}} & không & DF 'restricted' & open: dùng tự do; restricted: cần giấy phép \\
\texttt{description} & {\footnotesize\texttt{nvarchar(1000)}} & không & DF N'' & Mô tả, nguồn gốc \\
\texttt{created\_at} & {\footnotesize\texttt{datetime2(0)}} & không & DF sysdatetime() & Thời điểm tạo \\
\end{longtable}

\textbf{Bảng \texttt{topics}} -- Cây chủ đề phân loại video (quan hệ đệ quy qua parent\_topic\_id).

\setlength{\tabcolsep}{4pt}
\begin{longtable}{>{\raggedright\arraybackslash}p{0.21\textwidth}>{\raggedright\arraybackslash}p{0.19\textwidth}>{\centering\arraybackslash}p{0.07\textwidth}>{\raggedright\arraybackslash}p{0.17\textwidth}>{\raggedright\arraybackslash}p{0.26\textwidth}}
\caption{Cấu trúc bảng \texttt{topics}}\label{tab:col_topics}\\
\toprule \textbf{Cột} & \textbf{Kiểu} & \textbf{Null} & \textbf{Khoá / mặc định} & \textbf{Mô tả} \\ \midrule \endfirsthead
\toprule \textbf{Cột} & \textbf{Kiểu} & \textbf{Null} & \textbf{Khoá / mặc định} & \textbf{Mô tả} \\ \midrule \endhead
\bottomrule \endfoot
\texttt{topic\_id} & {\footnotesize\texttt{smallint}} & không & PK, IDENTITY & Mã chủ đề (tự tăng) \\
\texttt{topic\_name} & {\footnotesize\texttt{nvarchar(100)}} & không & UQ & Tên chủ đề \\
\texttt{parent\_topic\_id} & {\footnotesize\texttt{smallint}} & có & FK $\rightarrow$ topics & Chủ đề cha (NULL nếu là gốc) \\
\end{longtable}

\textbf{Bảng \texttt{video\_topics}} -- Quan hệ nhiều--nhiều giữa video và chủ đề.

\setlength{\tabcolsep}{4pt}
\begin{longtable}{>{\raggedright\arraybackslash}p{0.21\textwidth}>{\raggedright\arraybackslash}p{0.19\textwidth}>{\centering\arraybackslash}p{0.07\textwidth}>{\raggedright\arraybackslash}p{0.17\textwidth}>{\raggedright\arraybackslash}p{0.26\textwidth}}
\caption{Cấu trúc bảng \texttt{video\_topics}}\label{tab:col_video_topics}\\
\toprule \textbf{Cột} & \textbf{Kiểu} & \textbf{Null} & \textbf{Khoá / mặc định} & \textbf{Mô tả} \\ \midrule \endfirsthead
\toprule \textbf{Cột} & \textbf{Kiểu} & \textbf{Null} & \textbf{Khoá / mặc định} & \textbf{Mô tả} \\ \midrule \endhead
\bottomrule \endfoot
\texttt{video\_id} & {\footnotesize\texttt{varchar(16)}} & không & PK, FK $\rightarrow$ videos & Video \\
\texttt{topic\_id} & {\footnotesize\texttt{smallint}} & không & PK, FK $\rightarrow$ topics & Chủ đề gắn cho video \\
\end{longtable}

\subsection{Nhóm kho dữ liệu video}

\textbf{Bảng \texttt{videos}} -- Video trong kho, thuộc một bộ sưu tập.

\setlength{\tabcolsep}{4pt}
\begin{longtable}{>{\raggedright\arraybackslash}p{0.21\textwidth}>{\raggedright\arraybackslash}p{0.19\textwidth}>{\centering\arraybackslash}p{0.07\textwidth}>{\raggedright\arraybackslash}p{0.17\textwidth}>{\raggedright\arraybackslash}p{0.26\textwidth}}
\caption{Cấu trúc bảng \texttt{videos}}\label{tab:col_videos}\\
\toprule \textbf{Cột} & \textbf{Kiểu} & \textbf{Null} & \textbf{Khoá / mặc định} & \textbf{Mô tả} \\ \midrule \endfirsthead
\toprule \textbf{Cột} & \textbf{Kiểu} & \textbf{Null} & \textbf{Khoá / mặc định} & \textbf{Mô tả} \\ \midrule \endhead
\bottomrule \endfoot
\texttt{video\_id} & {\footnotesize\texttt{varchar(16)}} & không & PK & Mã video, dạng Lxx\_Vyyy \\
\texttt{batch} & {\footnotesize\texttt{tính toán}} & có &  & Lô dữ liệu (L23...), cột tính toán từ video\_id \\
\texttt{collection\_id} & {\footnotesize\texttt{int}} & không & FK $\rightarrow$ collections & Bộ sưu tập chứa video \\
\texttt{title} & {\footnotesize\texttt{nvarchar(200)}} & không &  & Nhan đề trong phiếu mô tả \\
\texttt{fps} & {\footnotesize\texttt{decimal(6,3)}} & có &  & Số khung hình mỗi giây của video gốc \\
\end{longtable}

\textbf{Bảng \texttt{keyframes}} -- Khung hình được trích từ video; khoá chính trùng số thứ tự vector trong FAISS.

\setlength{\tabcolsep}{4pt}
\begin{longtable}{>{\raggedright\arraybackslash}p{0.21\textwidth}>{\raggedright\arraybackslash}p{0.19\textwidth}>{\centering\arraybackslash}p{0.07\textwidth}>{\raggedright\arraybackslash}p{0.17\textwidth}>{\raggedright\arraybackslash}p{0.26\textwidth}}
\caption{Cấu trúc bảng \texttt{keyframes}}\label{tab:col_keyframes}\\
\toprule \textbf{Cột} & \textbf{Kiểu} & \textbf{Null} & \textbf{Khoá / mặc định} & \textbf{Mô tả} \\ \midrule \endfirsthead
\toprule \textbf{Cột} & \textbf{Kiểu} & \textbf{Null} & \textbf{Khoá / mặc định} & \textbf{Mô tả} \\ \midrule \endhead
\bottomrule \endfoot
\texttt{keyframe\_id} & {\footnotesize\texttt{int}} & không & PK & Mã khung hình, trùng id vector FAISS \\
\texttt{video\_id} & {\footnotesize\texttt{varchar(16)}} & không & UQ, FK $\rightarrow$ videos & Video chứa khung hình \\
\texttt{sample\_n} & {\footnotesize\texttt{int}} & không & UQ & Số thứ tự khung hình trong video (tên file ảnh) \\
\texttt{pts\_time} & {\footnotesize\texttt{decimal(10,3)}} & không & UQ & Thời điểm trong video (giây) \\
\texttt{frame\_idx} & {\footnotesize\texttt{int}} & không & UQ & Chỉ số khung hình trong video gốc (dùng khi trích đoạn) \\
\end{longtable}

\textbf{Bảng \texttt{ocr\_texts}} -- Chữ nhận dạng được trên khung hình (OCR).

\setlength{\tabcolsep}{4pt}
\begin{longtable}{>{\raggedright\arraybackslash}p{0.21\textwidth}>{\raggedright\arraybackslash}p{0.19\textwidth}>{\centering\arraybackslash}p{0.07\textwidth}>{\raggedright\arraybackslash}p{0.17\textwidth}>{\raggedright\arraybackslash}p{0.26\textwidth}}
\caption{Cấu trúc bảng \texttt{ocr\_texts}}\label{tab:col_ocr_texts}\\
\toprule \textbf{Cột} & \textbf{Kiểu} & \textbf{Null} & \textbf{Khoá / mặc định} & \textbf{Mô tả} \\ \midrule \endfirsthead
\toprule \textbf{Cột} & \textbf{Kiểu} & \textbf{Null} & \textbf{Khoá / mặc định} & \textbf{Mô tả} \\ \midrule \endhead
\bottomrule \endfoot
\texttt{keyframe\_id} & {\footnotesize\texttt{int}} & không & PK, FK $\rightarrow$ keyframes & Khung hình chứa chữ \\
\texttt{ocr\_text} & {\footnotesize\texttt{nvarchar(2000)}} & không &  & Chữ nhận dạng được \\
\texttt{text\_folded} & {\footnotesize\texttt{tính toán}} & có &  & Chữ đã bỏ dấu, viết thường (cột tính toán) \\
\end{longtable}

\textbf{Bảng \texttt{asr\_segments}} -- Đoạn lời nói nhận dạng được (ASR) theo khoảng thời gian.

\setlength{\tabcolsep}{4pt}
\begin{longtable}{>{\raggedright\arraybackslash}p{0.21\textwidth}>{\raggedright\arraybackslash}p{0.19\textwidth}>{\centering\arraybackslash}p{0.07\textwidth}>{\raggedright\arraybackslash}p{0.17\textwidth}>{\raggedright\arraybackslash}p{0.26\textwidth}}
\caption{Cấu trúc bảng \texttt{asr\_segments}}\label{tab:col_asr_segments}\\
\toprule \textbf{Cột} & \textbf{Kiểu} & \textbf{Null} & \textbf{Khoá / mặc định} & \textbf{Mô tả} \\ \midrule \endfirsthead
\toprule \textbf{Cột} & \textbf{Kiểu} & \textbf{Null} & \textbf{Khoá / mặc định} & \textbf{Mô tả} \\ \midrule \endhead
\bottomrule \endfoot
\texttt{segment\_id} & {\footnotesize\texttt{int}} & không & PK, IDENTITY & Mã đoạn lời nói (tự tăng) \\
\texttt{video\_id} & {\footnotesize\texttt{varchar(16)}} & không & FK $\rightarrow$ videos & Video chứa đoạn lời nói \\
\texttt{start\_sec} & {\footnotesize\texttt{decimal(10,3)}} & không &  & Thời điểm bắt đầu (giây) \\
\texttt{end\_sec} & {\footnotesize\texttt{decimal(10,3)}} & không &  & Thời điểm kết thúc (giây) \\
\texttt{asr\_text} & {\footnotesize\texttt{nvarchar(2000)}} & không &  & Lời nói nhận dạng được \\
\texttt{text\_folded} & {\footnotesize\texttt{tính toán}} & có &  & Lời nói đã bỏ dấu (cột tính toán) \\
\end{longtable}

\subsection{Nhóm chỉ mục văn bản}

\textbf{Bảng \texttt{stopwords}} -- Hư từ bị bỏ qua khi tách từ.

\setlength{\tabcolsep}{4pt}
\begin{longtable}{>{\raggedright\arraybackslash}p{0.21\textwidth}>{\raggedright\arraybackslash}p{0.19\textwidth}>{\centering\arraybackslash}p{0.07\textwidth}>{\raggedright\arraybackslash}p{0.17\textwidth}>{\raggedright\arraybackslash}p{0.26\textwidth}}
\caption{Cấu trúc bảng \texttt{stopwords}}\label{tab:col_stopwords}\\
\toprule \textbf{Cột} & \textbf{Kiểu} & \textbf{Null} & \textbf{Khoá / mặc định} & \textbf{Mô tả} \\ \midrule \endfirsthead
\toprule \textbf{Cột} & \textbf{Kiểu} & \textbf{Null} & \textbf{Khoá / mặc định} & \textbf{Mô tả} \\ \midrule \endhead
\bottomrule \endfoot
\texttt{word} & {\footnotesize\texttt{nvarchar(32)}} & không & PK & Hư từ đã bỏ dấu \\
\end{longtable}

\textbf{Bảng \texttt{terms}} -- Từ điển các từ khoá đã bỏ dấu.

\setlength{\tabcolsep}{4pt}
\begin{longtable}{>{\raggedright\arraybackslash}p{0.21\textwidth}>{\raggedright\arraybackslash}p{0.19\textwidth}>{\centering\arraybackslash}p{0.07\textwidth}>{\raggedright\arraybackslash}p{0.17\textwidth}>{\raggedright\arraybackslash}p{0.26\textwidth}}
\caption{Cấu trúc bảng \texttt{terms}}\label{tab:col_terms}\\
\toprule \textbf{Cột} & \textbf{Kiểu} & \textbf{Null} & \textbf{Khoá / mặc định} & \textbf{Mô tả} \\ \midrule \endfirsthead
\toprule \textbf{Cột} & \textbf{Kiểu} & \textbf{Null} & \textbf{Khoá / mặc định} & \textbf{Mô tả} \\ \midrule \endhead
\bottomrule \endfoot
\texttt{term\_id} & {\footnotesize\texttt{int}} & không & PK, IDENTITY & Mã từ khoá (tự tăng) \\
\texttt{term} & {\footnotesize\texttt{nvarchar(64)}} & không & UQ & Từ khoá đã bỏ dấu, so khớp nhị phân \\
\end{longtable}

\textbf{Bảng \texttt{ocr\_postings}} -- Chỉ mục ngược: từ khoá xuất hiện ở khung hình nào.

\setlength{\tabcolsep}{4pt}
\begin{longtable}{>{\raggedright\arraybackslash}p{0.21\textwidth}>{\raggedright\arraybackslash}p{0.19\textwidth}>{\centering\arraybackslash}p{0.07\textwidth}>{\raggedright\arraybackslash}p{0.17\textwidth}>{\raggedright\arraybackslash}p{0.26\textwidth}}
\caption{Cấu trúc bảng \texttt{ocr\_postings}}\label{tab:col_ocr_postings}\\
\toprule \textbf{Cột} & \textbf{Kiểu} & \textbf{Null} & \textbf{Khoá / mặc định} & \textbf{Mô tả} \\ \midrule \endfirsthead
\toprule \textbf{Cột} & \textbf{Kiểu} & \textbf{Null} & \textbf{Khoá / mặc định} & \textbf{Mô tả} \\ \midrule \endhead
\bottomrule \endfoot
\texttt{term\_id} & {\footnotesize\texttt{int}} & không & PK, FK $\rightarrow$ terms & Từ khoá \\
\texttt{keyframe\_id} & {\footnotesize\texttt{int}} & không & PK, FK $\rightarrow$ ocr\_texts & Khung hình có từ khoá trong chữ OCR \\
\end{longtable}

\textbf{Bảng \texttt{asr\_postings}} -- Chỉ mục ngược: từ khoá xuất hiện ở đoạn lời nói nào.

\setlength{\tabcolsep}{4pt}
\begin{longtable}{>{\raggedright\arraybackslash}p{0.21\textwidth}>{\raggedright\arraybackslash}p{0.19\textwidth}>{\centering\arraybackslash}p{0.07\textwidth}>{\raggedright\arraybackslash}p{0.17\textwidth}>{\raggedright\arraybackslash}p{0.26\textwidth}}
\caption{Cấu trúc bảng \texttt{asr\_postings}}\label{tab:col_asr_postings}\\
\toprule \textbf{Cột} & \textbf{Kiểu} & \textbf{Null} & \textbf{Khoá / mặc định} & \textbf{Mô tả} \\ \midrule \endfirsthead
\toprule \textbf{Cột} & \textbf{Kiểu} & \textbf{Null} & \textbf{Khoá / mặc định} & \textbf{Mô tả} \\ \midrule \endhead
\bottomrule \endfoot
\texttt{term\_id} & {\footnotesize\texttt{int}} & không & PK, FK $\rightarrow$ terms & Từ khoá \\
\texttt{segment\_id} & {\footnotesize\texttt{int}} & không & PK, FK $\rightarrow$ asr\_segments & Đoạn lời nói có từ khoá \\
\end{longtable}

\subsection{Nhóm tài khoản và phân quyền}

\textbf{Bảng \texttt{users}} -- Tài khoản đăng nhập: bạn đọc và cán bộ thư viện.

\setlength{\tabcolsep}{4pt}
\begin{longtable}{>{\raggedright\arraybackslash}p{0.21\textwidth}>{\raggedright\arraybackslash}p{0.19\textwidth}>{\centering\arraybackslash}p{0.07\textwidth}>{\raggedright\arraybackslash}p{0.17\textwidth}>{\raggedright\arraybackslash}p{0.26\textwidth}}
\caption{Cấu trúc bảng \texttt{users}}\label{tab:col_users}\\
\toprule \textbf{Cột} & \textbf{Kiểu} & \textbf{Null} & \textbf{Khoá / mặc định} & \textbf{Mô tả} \\ \midrule \endfirsthead
\toprule \textbf{Cột} & \textbf{Kiểu} & \textbf{Null} & \textbf{Khoá / mặc định} & \textbf{Mô tả} \\ \midrule \endhead
\bottomrule \endfoot
\texttt{user\_id} & {\footnotesize\texttt{int}} & không & PK, IDENTITY & Mã người dùng (tự tăng) \\
\texttt{username} & {\footnotesize\texttt{varchar(32)}} & không & UQ & Tên đăng nhập: 3--32 ký tự a-z, 0-9, \_ . - \\
\texttt{full\_name} & {\footnotesize\texttt{nvarchar(100)}} & không & DF N'' & Họ tên hiển thị \\
\texttt{password\_hash} & {\footnotesize\texttt{varchar(255)}} & không &  & Mật khẩu băm scrypt (kèm muối) \\
\texttt{organization\_id} & {\footnotesize\texttt{smallint}} & có & FK $\rightarrow$ organizations & Đơn vị công tác (NULL: bạn đọc cá nhân) \\
\texttt{user\_type} & {\footnotesize\texttt{varchar(16)}} & không & DF 'other' & Đối tượng: staff, journalist, researcher, student, other \\
\texttt{is\_active} & {\footnotesize\texttt{bit}} & không & DF 1 & Tài khoản còn hoạt động \\
\texttt{created\_at} & {\footnotesize\texttt{datetime2(0)}} & không & DF sysdatetime() & Thời điểm tạo \\
\texttt{last\_login\_at} & {\footnotesize\texttt{datetime2(0)}} & có &  & Lần đăng nhập gần nhất \\
\end{longtable}

\textbf{Bảng \texttt{roles}} -- Vai trò (reader, librarian, admin).

\setlength{\tabcolsep}{4pt}
\begin{longtable}{>{\raggedright\arraybackslash}p{0.21\textwidth}>{\raggedright\arraybackslash}p{0.19\textwidth}>{\centering\arraybackslash}p{0.07\textwidth}>{\raggedright\arraybackslash}p{0.17\textwidth}>{\raggedright\arraybackslash}p{0.26\textwidth}}
\caption{Cấu trúc bảng \texttt{roles}}\label{tab:col_roles}\\
\toprule \textbf{Cột} & \textbf{Kiểu} & \textbf{Null} & \textbf{Khoá / mặc định} & \textbf{Mô tả} \\ \midrule \endfirsthead
\toprule \textbf{Cột} & \textbf{Kiểu} & \textbf{Null} & \textbf{Khoá / mặc định} & \textbf{Mô tả} \\ \midrule \endhead
\bottomrule \endfoot
\texttt{role\_id} & {\footnotesize\texttt{tinyint}} & không & PK, IDENTITY & Mã vai trò (tự tăng) \\
\texttt{role\_name} & {\footnotesize\texttt{varchar(32)}} & không & UQ & Tên vai trò: reader, librarian, admin \\
\texttt{description} & {\footnotesize\texttt{nvarchar(255)}} & không & DF N'' & Mô tả \\
\end{longtable}

\textbf{Bảng \texttt{permissions}} -- Quyền nghiệp vụ.

\setlength{\tabcolsep}{4pt}
\begin{longtable}{>{\raggedright\arraybackslash}p{0.21\textwidth}>{\raggedright\arraybackslash}p{0.19\textwidth}>{\centering\arraybackslash}p{0.07\textwidth}>{\raggedright\arraybackslash}p{0.17\textwidth}>{\raggedright\arraybackslash}p{0.26\textwidth}}
\caption{Cấu trúc bảng \texttt{permissions}}\label{tab:col_permissions}\\
\toprule \textbf{Cột} & \textbf{Kiểu} & \textbf{Null} & \textbf{Khoá / mặc định} & \textbf{Mô tả} \\ \midrule \endfirsthead
\toprule \textbf{Cột} & \textbf{Kiểu} & \textbf{Null} & \textbf{Khoá / mặc định} & \textbf{Mô tả} \\ \midrule \endhead
\bottomrule \endfoot
\texttt{permission\_code} & {\footnotesize\texttt{varchar(32)}} & không & PK & Mã quyền, vd. selection.write \\
\texttt{description} & {\footnotesize\texttt{nvarchar(255)}} & không & DF N'' & Mô tả \\
\end{longtable}

\textbf{Bảng \texttt{user\_roles}} -- Quan hệ nhiều--nhiều giữa người dùng và vai trò.

\setlength{\tabcolsep}{4pt}
\begin{longtable}{>{\raggedright\arraybackslash}p{0.21\textwidth}>{\raggedright\arraybackslash}p{0.19\textwidth}>{\centering\arraybackslash}p{0.07\textwidth}>{\raggedright\arraybackslash}p{0.17\textwidth}>{\raggedright\arraybackslash}p{0.26\textwidth}}
\caption{Cấu trúc bảng \texttt{user\_roles}}\label{tab:col_user_roles}\\
\toprule \textbf{Cột} & \textbf{Kiểu} & \textbf{Null} & \textbf{Khoá / mặc định} & \textbf{Mô tả} \\ \midrule \endfirsthead
\toprule \textbf{Cột} & \textbf{Kiểu} & \textbf{Null} & \textbf{Khoá / mặc định} & \textbf{Mô tả} \\ \midrule \endhead
\bottomrule \endfoot
\texttt{user\_id} & {\footnotesize\texttt{int}} & không & PK, FK $\rightarrow$ users & Người dùng \\
\texttt{role\_id} & {\footnotesize\texttt{tinyint}} & không & PK, FK $\rightarrow$ roles & Vai trò được gán \\
\end{longtable}

\textbf{Bảng \texttt{role\_permissions}} -- Quan hệ nhiều--nhiều giữa vai trò và quyền.

\setlength{\tabcolsep}{4pt}
\begin{longtable}{>{\raggedright\arraybackslash}p{0.21\textwidth}>{\raggedright\arraybackslash}p{0.19\textwidth}>{\centering\arraybackslash}p{0.07\textwidth}>{\raggedright\arraybackslash}p{0.17\textwidth}>{\raggedright\arraybackslash}p{0.26\textwidth}}
\caption{Cấu trúc bảng \texttt{role\_permissions}}\label{tab:col_role_permissions}\\
\toprule \textbf{Cột} & \textbf{Kiểu} & \textbf{Null} & \textbf{Khoá / mặc định} & \textbf{Mô tả} \\ \midrule \endfirsthead
\toprule \textbf{Cột} & \textbf{Kiểu} & \textbf{Null} & \textbf{Khoá / mặc định} & \textbf{Mô tả} \\ \midrule \endhead
\bottomrule \endfoot
\texttt{role\_id} & {\footnotesize\texttt{tinyint}} & không & PK, FK $\rightarrow$ roles & Vai trò \\
\texttt{permission\_code} & {\footnotesize\texttt{varchar(32)}} & không & PK, FK $\rightarrow$ permissions & Quyền của vai trò \\
\end{longtable}

\textbf{Bảng \texttt{sessions}} -- Phiên đăng nhập (chỉ lưu băm của cookie).

\setlength{\tabcolsep}{4pt}
\begin{longtable}{>{\raggedright\arraybackslash}p{0.21\textwidth}>{\raggedright\arraybackslash}p{0.19\textwidth}>{\centering\arraybackslash}p{0.07\textwidth}>{\raggedright\arraybackslash}p{0.17\textwidth}>{\raggedright\arraybackslash}p{0.26\textwidth}}
\caption{Cấu trúc bảng \texttt{sessions}}\label{tab:col_sessions}\\
\toprule \textbf{Cột} & \textbf{Kiểu} & \textbf{Null} & \textbf{Khoá / mặc định} & \textbf{Mô tả} \\ \midrule \endfirsthead
\toprule \textbf{Cột} & \textbf{Kiểu} & \textbf{Null} & \textbf{Khoá / mặc định} & \textbf{Mô tả} \\ \midrule \endhead
\bottomrule \endfoot
\texttt{session\_hash} & {\footnotesize\texttt{char(64)}} & không & PK & SHA-256 của cookie đăng nhập \\
\texttt{user\_id} & {\footnotesize\texttt{int}} & không & FK $\rightarrow$ users & Chủ phiên \\
\texttt{created\_at} & {\footnotesize\texttt{datetime2(0)}} & không & DF sysdatetime() & Thời điểm đăng nhập \\
\texttt{expires\_at} & {\footnotesize\texttt{datetime2(0)}} & không &  & Thời điểm hết hạn \\
\end{longtable}

\subsection{Nhóm yêu cầu khai thác và cấp phép}

\textbf{Bảng \texttt{requests}} -- Yêu cầu khai thác: bạn đọc mô tả đoạn tư liệu cần dùng, mục đích và trạng thái xử lý.

\setlength{\tabcolsep}{4pt}
\begin{longtable}{>{\raggedright\arraybackslash}p{0.21\textwidth}>{\raggedright\arraybackslash}p{0.19\textwidth}>{\centering\arraybackslash}p{0.07\textwidth}>{\raggedright\arraybackslash}p{0.17\textwidth}>{\raggedright\arraybackslash}p{0.26\textwidth}}
\caption{Cấu trúc bảng \texttt{requests}}\label{tab:col_requests}\\
\toprule \textbf{Cột} & \textbf{Kiểu} & \textbf{Null} & \textbf{Khoá / mặc định} & \textbf{Mô tả} \\ \midrule \endfirsthead
\toprule \textbf{Cột} & \textbf{Kiểu} & \textbf{Null} & \textbf{Khoá / mặc định} & \textbf{Mô tả} \\ \midrule \endhead
\bottomrule \endfoot
\texttt{request\_id} & {\footnotesize\texttt{varchar(32)}} & không & PK & Mã yêu cầu, vd. YC-047, YC-U0001 \\
\texttt{request\_text} & {\footnotesize\texttt{nvarchar(2000)}} & không &  & Mô tả đoạn tư liệu cần dùng \\
\texttt{origin} & {\footnotesize\texttt{varchar(16)}} & không &  & Nguồn: sample (mẫu đã xác minh) hoặc user (gửi qua web) \\
\texttt{requested\_by} & {\footnotesize\texttt{int}} & không & FK $\rightarrow$ users & Bạn đọc gửi yêu cầu \\
\texttt{purpose} & {\footnotesize\texttt{varchar(16)}} & không & DF 'research' & Mục đích: broadcast, research, education, exhibition \\
\texttt{status} & {\footnotesize\texttt{varchar(16)}} & không & DF 'pending' & pending, approved, rejected; chỉ đổi qua thủ tục \\
\texttt{created\_at} & {\footnotesize\texttt{datetime2(0)}} & không & DF sysdatetime() & Thời điểm gửi \\
\end{longtable}

\textbf{Bảng \texttt{ground\_truth}} -- Vị trí đúng (xác minh thủ công) của đoạn tư liệu mà yêu cầu mẫu cần tìm.

\setlength{\tabcolsep}{4pt}
\begin{longtable}{>{\raggedright\arraybackslash}p{0.21\textwidth}>{\raggedright\arraybackslash}p{0.19\textwidth}>{\centering\arraybackslash}p{0.07\textwidth}>{\raggedright\arraybackslash}p{0.17\textwidth}>{\raggedright\arraybackslash}p{0.26\textwidth}}
\caption{Cấu trúc bảng \texttt{ground\_truth}}\label{tab:col_ground_truth}\\
\toprule \textbf{Cột} & \textbf{Kiểu} & \textbf{Null} & \textbf{Khoá / mặc định} & \textbf{Mô tả} \\ \midrule \endfirsthead
\toprule \textbf{Cột} & \textbf{Kiểu} & \textbf{Null} & \textbf{Khoá / mặc định} & \textbf{Mô tả} \\ \midrule \endhead
\bottomrule \endfoot
\texttt{request\_id} & {\footnotesize\texttt{varchar(32)}} & không & PK, FK $\rightarrow$ requests & Yêu cầu mẫu \\
\texttt{video\_id} & {\footnotesize\texttt{varchar(16)}} & không & FK $\rightarrow$ videos & Video chứa cảnh đúng \\
\texttt{start\_sec} & {\footnotesize\texttt{decimal(10,3)}} & không &  & Bắt đầu khoảng thời gian đúng (giây) \\
\texttt{end\_sec} & {\footnotesize\texttt{decimal(10,3)}} & không &  & Kết thúc khoảng thời gian đúng (giây) \\
\end{longtable}

\textbf{Bảng \texttt{selections}} -- Đoạn tư liệu (khung hình) cán bộ thư viện chọn cho một yêu cầu.

\setlength{\tabcolsep}{4pt}
\begin{longtable}{>{\raggedright\arraybackslash}p{0.21\textwidth}>{\raggedright\arraybackslash}p{0.19\textwidth}>{\centering\arraybackslash}p{0.07\textwidth}>{\raggedright\arraybackslash}p{0.17\textwidth}>{\raggedright\arraybackslash}p{0.26\textwidth}}
\caption{Cấu trúc bảng \texttt{selections}}\label{tab:col_selections}\\
\toprule \textbf{Cột} & \textbf{Kiểu} & \textbf{Null} & \textbf{Khoá / mặc định} & \textbf{Mô tả} \\ \midrule \endfirsthead
\toprule \textbf{Cột} & \textbf{Kiểu} & \textbf{Null} & \textbf{Khoá / mặc định} & \textbf{Mô tả} \\ \midrule \endhead
\bottomrule \endfoot
\texttt{selection\_id} & {\footnotesize\texttt{int}} & không & PK, IDENTITY & Mã lựa chọn (tự tăng) \\
\texttt{request\_id} & {\footnotesize\texttt{varchar(32)}} & không & UQ, FK $\rightarrow$ requests & Yêu cầu khai thác \\
\texttt{keyframe\_id} & {\footnotesize\texttt{int}} & không & UQ, FK $\rightarrow$ keyframes & Đoạn tư liệu (khung hình) được chọn \\
\texttt{user\_id} & {\footnotesize\texttt{int}} & có & FK $\rightarrow$ users & Người chọn (NULL nếu tài khoản đã bị xoá) \\
\texttt{note} & {\footnotesize\texttt{nvarchar(255)}} & không & DF N'' & Ghi chú \\
\texttt{created\_at} & {\footnotesize\texttt{datetime2(0)}} & không & DF sysdatetime() & Thời điểm chọn; quyết định thứ tự trong danh sách xuất \\
\end{longtable}

\textbf{Bảng \texttt{usage\_licenses}} -- Giấy phép sử dụng cấp cho một yêu cầu đã duyệt.

\setlength{\tabcolsep}{4pt}
\begin{longtable}{>{\raggedright\arraybackslash}p{0.21\textwidth}>{\raggedright\arraybackslash}p{0.19\textwidth}>{\centering\arraybackslash}p{0.07\textwidth}>{\raggedright\arraybackslash}p{0.17\textwidth}>{\raggedright\arraybackslash}p{0.26\textwidth}}
\caption{Cấu trúc bảng \texttt{usage\_licenses}}\label{tab:col_usage_licenses}\\
\toprule \textbf{Cột} & \textbf{Kiểu} & \textbf{Null} & \textbf{Khoá / mặc định} & \textbf{Mô tả} \\ \midrule \endfirsthead
\toprule \textbf{Cột} & \textbf{Kiểu} & \textbf{Null} & \textbf{Khoá / mặc định} & \textbf{Mô tả} \\ \midrule \endhead
\bottomrule \endfoot
\texttt{license\_id} & {\footnotesize\texttt{int}} & không & PK, IDENTITY & Mã giấy phép (tự tăng) \\
\texttt{request\_id} & {\footnotesize\texttt{varchar(32)}} & không & UQ, FK $\rightarrow$ requests & Yêu cầu được cấp phép (duy nhất) \\
\texttt{licensee\_org\_id} & {\footnotesize\texttt{smallint}} & có & FK $\rightarrow$ organizations & Đơn vị được cấp phép tại thời điểm cấp (NULL: cá nhân) \\
\texttt{granted\_by} & {\footnotesize\texttt{int}} & không & FK $\rightarrow$ users & Cán bộ cấp phép \\
\texttt{granted\_at} & {\footnotesize\texttt{datetime2(0)}} & không & DF sysdatetime() & Thời điểm cấp \\
\texttt{valid\_from} & {\footnotesize\texttt{date}} & không &  & Ngày bắt đầu hiệu lực \\
\texttt{valid\_to} & {\footnotesize\texttt{date}} & không &  & Ngày hết hiệu lực (sau valid\_from) \\
\texttt{conditions} & {\footnotesize\texttt{nvarchar(500)}} & không & DF N'' & Điều kiện sử dụng \\
\end{longtable}

\subsection{Nhóm nhật ký và đánh giá}

\textbf{Bảng \texttt{sources}} -- Danh mục nguồn tìm kiếm (visual, ocr, asr).

\setlength{\tabcolsep}{4pt}
\begin{longtable}{>{\raggedright\arraybackslash}p{0.21\textwidth}>{\raggedright\arraybackslash}p{0.19\textwidth}>{\centering\arraybackslash}p{0.07\textwidth}>{\raggedright\arraybackslash}p{0.17\textwidth}>{\raggedright\arraybackslash}p{0.26\textwidth}}
\caption{Cấu trúc bảng \texttt{sources}}\label{tab:col_sources}\\
\toprule \textbf{Cột} & \textbf{Kiểu} & \textbf{Null} & \textbf{Khoá / mặc định} & \textbf{Mô tả} \\ \midrule \endfirsthead
\toprule \textbf{Cột} & \textbf{Kiểu} & \textbf{Null} & \textbf{Khoá / mặc định} & \textbf{Mô tả} \\ \midrule \endhead
\bottomrule \endfoot
\texttt{source\_code} & {\footnotesize\texttt{varchar(16)}} & không & PK & Mã nguồn: visual, ocr, asr \\
\texttt{description} & {\footnotesize\texttt{nvarchar(100)}} & không & DF N'' & Mô tả \\
\end{longtable}

\textbf{Bảng \texttt{audit\_log}} -- Nhật ký thao tác, do trigger và ứng dụng ghi.

\setlength{\tabcolsep}{4pt}
\begin{longtable}{>{\raggedright\arraybackslash}p{0.21\textwidth}>{\raggedright\arraybackslash}p{0.19\textwidth}>{\centering\arraybackslash}p{0.07\textwidth}>{\raggedright\arraybackslash}p{0.17\textwidth}>{\raggedright\arraybackslash}p{0.26\textwidth}}
\caption{Cấu trúc bảng \texttt{audit\_log}}\label{tab:col_audit_log}\\
\toprule \textbf{Cột} & \textbf{Kiểu} & \textbf{Null} & \textbf{Khoá / mặc định} & \textbf{Mô tả} \\ \midrule \endfirsthead
\toprule \textbf{Cột} & \textbf{Kiểu} & \textbf{Null} & \textbf{Khoá / mặc định} & \textbf{Mô tả} \\ \midrule \endhead
\bottomrule \endfoot
\texttt{log\_id} & {\footnotesize\texttt{bigint}} & không & PK, IDENTITY & Mã dòng nhật ký (tự tăng) \\
\texttt{user\_id} & {\footnotesize\texttt{int}} & có & FK $\rightarrow$ users & Người thực hiện \\
\texttt{action} & {\footnotesize\texttt{varchar(32)}} & không &  & Hành động, vd. selection.add, request.approve \\
\texttt{detail} & {\footnotesize\texttt{nvarchar(500)}} & không & DF N'' & Chi tiết dạng văn bản \\
\texttt{created\_at} & {\footnotesize\texttt{datetime2(0)}} & không & DF sysdatetime() & Thời điểm \\
\end{longtable}

\textbf{Bảng \texttt{search\_log}} -- Nhật ký các lượt tìm kiếm và độ trễ.

\setlength{\tabcolsep}{4pt}
\begin{longtable}{>{\raggedright\arraybackslash}p{0.21\textwidth}>{\raggedright\arraybackslash}p{0.19\textwidth}>{\centering\arraybackslash}p{0.07\textwidth}>{\raggedright\arraybackslash}p{0.17\textwidth}>{\raggedright\arraybackslash}p{0.26\textwidth}}
\caption{Cấu trúc bảng \texttt{search\_log}}\label{tab:col_search_log}\\
\toprule \textbf{Cột} & \textbf{Kiểu} & \textbf{Null} & \textbf{Khoá / mặc định} & \textbf{Mô tả} \\ \midrule \endfirsthead
\toprule \textbf{Cột} & \textbf{Kiểu} & \textbf{Null} & \textbf{Khoá / mặc định} & \textbf{Mô tả} \\ \midrule \endhead
\bottomrule \endfoot
\texttt{search\_id} & {\footnotesize\texttt{bigint}} & không & PK, IDENTITY & Mã lượt tìm (tự tăng) \\
\texttt{user\_id} & {\footnotesize\texttt{int}} & có & FK $\rightarrow$ users & Người tìm \\
\texttt{request\_text} & {\footnotesize\texttt{nvarchar(2000)}} & không &  & Nội dung tìm \\
\texttt{top\_k} & {\footnotesize\texttt{smallint}} & không &  & Số kết quả yêu cầu \\
\texttt{result\_count} & {\footnotesize\texttt{smallint}} & không &  & Số kết quả trả về \\
\texttt{latency\_ms} & {\footnotesize\texttt{int}} & không &  & Thời gian xử lý (ms) \\
\texttt{created\_at} & {\footnotesize\texttt{datetime2(0)}} & không & DF sysdatetime() & Thời điểm \\
\end{longtable}

\textbf{Bảng \texttt{search\_log\_sources}} -- Nguồn được dùng trong mỗi lượt tìm kiếm.

\setlength{\tabcolsep}{4pt}
\begin{longtable}{>{\raggedright\arraybackslash}p{0.21\textwidth}>{\raggedright\arraybackslash}p{0.19\textwidth}>{\centering\arraybackslash}p{0.07\textwidth}>{\raggedright\arraybackslash}p{0.17\textwidth}>{\raggedright\arraybackslash}p{0.26\textwidth}}
\caption{Cấu trúc bảng \texttt{search\_log\_sources}}\label{tab:col_search_log_sources}\\
\toprule \textbf{Cột} & \textbf{Kiểu} & \textbf{Null} & \textbf{Khoá / mặc định} & \textbf{Mô tả} \\ \midrule \endfirsthead
\toprule \textbf{Cột} & \textbf{Kiểu} & \textbf{Null} & \textbf{Khoá / mặc định} & \textbf{Mô tả} \\ \midrule \endhead
\bottomrule \endfoot
\texttt{search\_id} & {\footnotesize\texttt{bigint}} & không & PK, FK $\rightarrow$ search\_log & Lượt tìm \\
\texttt{source\_code} & {\footnotesize\texttt{varchar(16)}} & không & PK, FK $\rightarrow$ sources & Nguồn được dùng \\
\end{longtable}

\textbf{Bảng \texttt{eval\_runs}} -- Lượt đánh giá trên tập yêu cầu mẫu đã xác minh.

\setlength{\tabcolsep}{4pt}
\begin{longtable}{>{\raggedright\arraybackslash}p{0.21\textwidth}>{\raggedright\arraybackslash}p{0.19\textwidth}>{\centering\arraybackslash}p{0.07\textwidth}>{\raggedright\arraybackslash}p{0.17\textwidth}>{\raggedright\arraybackslash}p{0.26\textwidth}}
\caption{Cấu trúc bảng \texttt{eval\_runs}}\label{tab:col_eval_runs}\\
\toprule \textbf{Cột} & \textbf{Kiểu} & \textbf{Null} & \textbf{Khoá / mặc định} & \textbf{Mô tả} \\ \midrule \endfirsthead
\toprule \textbf{Cột} & \textbf{Kiểu} & \textbf{Null} & \textbf{Khoá / mặc định} & \textbf{Mô tả} \\ \midrule \endhead
\bottomrule \endfoot
\texttt{run\_id} & {\footnotesize\texttt{int}} & không & PK, IDENTITY & Mã lượt đánh giá (tự tăng) \\
\texttt{run\_name} & {\footnotesize\texttt{nvarchar(64)}} & không &  & Tên lượt \\
\texttt{top\_k} & {\footnotesize\texttt{smallint}} & không &  & Số kết quả xét cho mỗi yêu cầu \\
\texttt{tolerance\_sec} & {\footnotesize\texttt{decimal(5,2)}} & không & DF 3.0 & Sai lệch thời gian cho phép (giây) \\
\texttt{created\_at} & {\footnotesize\texttt{datetime2(0)}} & không & DF sysdatetime() & Thời điểm chạy \\
\end{longtable}

\textbf{Bảng \texttt{eval\_run\_sources}} -- Nguồn được dùng trong mỗi lượt đánh giá.

\setlength{\tabcolsep}{4pt}
\begin{longtable}{>{\raggedright\arraybackslash}p{0.21\textwidth}>{\raggedright\arraybackslash}p{0.19\textwidth}>{\centering\arraybackslash}p{0.07\textwidth}>{\raggedright\arraybackslash}p{0.17\textwidth}>{\raggedright\arraybackslash}p{0.26\textwidth}}
\caption{Cấu trúc bảng \texttt{eval\_run\_sources}}\label{tab:col_eval_run_sources}\\
\toprule \textbf{Cột} & \textbf{Kiểu} & \textbf{Null} & \textbf{Khoá / mặc định} & \textbf{Mô tả} \\ \midrule \endfirsthead
\toprule \textbf{Cột} & \textbf{Kiểu} & \textbf{Null} & \textbf{Khoá / mặc định} & \textbf{Mô tả} \\ \midrule \endhead
\bottomrule \endfoot
\texttt{run\_id} & {\footnotesize\texttt{int}} & không & PK, FK $\rightarrow$ eval\_runs & Lượt đánh giá \\
\texttt{source\_code} & {\footnotesize\texttt{varchar(16)}} & không & PK, FK $\rightarrow$ sources & Nguồn được dùng \\
\end{longtable}

\textbf{Bảng \texttt{eval\_requests}} -- Từng yêu cầu được chạy trong một lượt đánh giá và độ trễ.

\setlength{\tabcolsep}{4pt}
\begin{longtable}{>{\raggedright\arraybackslash}p{0.21\textwidth}>{\raggedright\arraybackslash}p{0.19\textwidth}>{\centering\arraybackslash}p{0.07\textwidth}>{\raggedright\arraybackslash}p{0.17\textwidth}>{\raggedright\arraybackslash}p{0.26\textwidth}}
\caption{Cấu trúc bảng \texttt{eval\_requests}}\label{tab:col_eval_requests}\\
\toprule \textbf{Cột} & \textbf{Kiểu} & \textbf{Null} & \textbf{Khoá / mặc định} & \textbf{Mô tả} \\ \midrule \endfirsthead
\toprule \textbf{Cột} & \textbf{Kiểu} & \textbf{Null} & \textbf{Khoá / mặc định} & \textbf{Mô tả} \\ \midrule \endhead
\bottomrule \endfoot
\texttt{run\_id} & {\footnotesize\texttt{int}} & không & PK, FK $\rightarrow$ eval\_runs & Lượt đánh giá \\
\texttt{request\_id} & {\footnotesize\texttt{varchar(32)}} & không & PK, FK $\rightarrow$ ground\_truth & Yêu cầu mẫu được chạy \\
\texttt{latency\_ms} & {\footnotesize\texttt{int}} & không &  & Thời gian xử lý (ms) \\
\end{longtable}

\textbf{Bảng \texttt{eval\_results}} -- Danh sách kết quả xếp hạng của từng yêu cầu trong lượt đánh giá.

\setlength{\tabcolsep}{4pt}
\begin{longtable}{>{\raggedright\arraybackslash}p{0.21\textwidth}>{\raggedright\arraybackslash}p{0.19\textwidth}>{\centering\arraybackslash}p{0.07\textwidth}>{\raggedright\arraybackslash}p{0.17\textwidth}>{\raggedright\arraybackslash}p{0.26\textwidth}}
\caption{Cấu trúc bảng \texttt{eval\_results}}\label{tab:col_eval_results}\\
\toprule \textbf{Cột} & \textbf{Kiểu} & \textbf{Null} & \textbf{Khoá / mặc định} & \textbf{Mô tả} \\ \midrule \endfirsthead
\toprule \textbf{Cột} & \textbf{Kiểu} & \textbf{Null} & \textbf{Khoá / mặc định} & \textbf{Mô tả} \\ \midrule \endhead
\bottomrule \endfoot
\texttt{run\_id} & {\footnotesize\texttt{int}} & không & PK, UQ, FK $\rightarrow$ eval\_requests & Lượt đánh giá \\
\texttt{request\_id} & {\footnotesize\texttt{varchar(32)}} & không & PK, UQ, FK $\rightarrow$ eval\_requests & Yêu cầu được chạy \\
\texttt{rank\_no} & {\footnotesize\texttt{smallint}} & không & PK & Thứ hạng của kết quả \\
\texttt{keyframe\_id} & {\footnotesize\texttt{int}} & không & UQ, FK $\rightarrow$ keyframes & Khung hình được trả về \\
\end{longtable}

}

\section{Ràng buộc toàn vẹn}

\subsection{Quy ước và số lượng}

Mọi ràng buộc đều được đặt tên theo quy ước \texttt{<loại>\_<bảng>\_<ý nghĩa>}:
\texttt{PK\_} khoá chính, \texttt{FK\_} khoá ngoại, \texttt{UQ\_} duy nhất, \texttt{CK\_}
kiểm tra, \texttt{DF\_} mặc định. Nhờ vậy thông báo lỗi chỉ rõ ràng buộc nào bị vi phạm
(xem mục~\ref{sec:tests}). Bảng~\ref{tab:constraint_counts} thống kê theo catalog.

\begin{table}[H]
\centering
\caption{Số lượng ràng buộc theo loại}
\label{tab:constraint_counts}
\begin{tabular}{lr}
\toprule \textbf{Loại ràng buộc} & \textbf{Số lượng} \\ \midrule
PRIMARY KEY (PK\_) & 30 \\
FOREIGN KEY (FK\_) & 37 \\
UNIQUE (UQ\_) & 13 \\
CHECK (CK\_) & 24 \\
DEFAULT (DF\_) & 24 \\
NOT NULL (cột) & 109 \\
\bottomrule
\end{tabular}
\end{table}

\subsection{Ví dụ khai báo}

Đoạn khai báo bảng \texttt{users} dùng đủ sáu loại ràng buộc ở mức cột và mức bảng:

\begin{lstlisting}[caption={Bảng users với đủ các loại ràng buộc}, label={lst:users}]
CREATE TABLE dbo.users (
    user_id         INT IDENTITY(1,1) NOT NULL,
    username        VARCHAR(32)   NOT NULL,
    full_name       NVARCHAR(100) NOT NULL CONSTRAINT DF_users_full_name DEFAULT N'',
    password_hash   VARCHAR(255)  NOT NULL,
    organization_id SMALLINT      NULL,
    user_type       VARCHAR(16)   NOT NULL CONSTRAINT DF_users_type DEFAULT 'other',
    is_active       BIT           NOT NULL CONSTRAINT DF_users_is_active DEFAULT 1,
    created_at      DATETIME2(0)  NOT NULL CONSTRAINT DF_users_created DEFAULT SYSDATETIME(),
    last_login_at   DATETIME2(0)  NULL,
    CONSTRAINT PK_users PRIMARY KEY (user_id),
    CONSTRAINT UQ_users_username UNIQUE (username),
    CONSTRAINT FK_users_organization FOREIGN KEY (organization_id)
        REFERENCES dbo.organizations (organization_id),
    CONSTRAINT CK_users_username CHECK (
        LEN(username) BETWEEN 3 AND 32
        AND username COLLATE Latin1_General_100_BIN2 NOT LIKE '%[^a-z0-9_.-]%'),
    CONSTRAINT CK_users_type CHECK (user_type IN ('staff', 'journalist', 'researcher', 'student', 'other'))
);
\end{lstlisting}

Một số điểm đáng chú ý:
\begin{itemize}
    \item \texttt{CK\_users\_username} so khớp với collation nhị phân để phân biệt hoa
    thường; collation mặc định không phân biệt nên \texttt{'Admin'} sẽ lọt qua.
    \item \texttt{CK\_videos\_id} dùng mẫu \texttt{'L[0-9][0-9][\_]V[0-9][0-9][0-9]'}, trong
    đó \texttt{[\_]} thoát ký tự \texttt{\_}, vốn là ký tự đại diện của \texttt{LIKE}.
    \item \texttt{CK\_usage\_licenses\_period} (\texttt{valid\_to > valid\_from}) là ràng buộc
    giữa hai cột nên khai báo ở mức bảng. \texttt{CK\_topics\_parent} chặn một chủ đề tự làm
    cha của chính nó.
    \item Hành vi khi xoá được chọn theo ngữ nghĩa. Xoá yêu cầu khai thác thì xoá luôn vị trí
    đã xác minh và các đoạn tư liệu đã chọn (\texttt{ON DELETE CASCADE}), nhưng không xoá được
    yêu cầu đã có giấy phép: giấy phép là chứng từ, khoá ngoại để mặc định \texttt{NO ACTION}.
    Không cho xoá khung hình đang được chọn cho một yêu cầu (\texttt{NO ACTION}).
    \item Các khoá ngoại tới \texttt{users} trong nhóm yêu cầu (\texttt{requests.requested\_by},
    \texttt{usage\_licenses.granted\_by}) đều \texttt{NO ACTION}: tài khoản được khoá chứ không
    bị xoá, và yêu cầu, giấy phép nào cũng phải có người gửi, người cấp (cột \texttt{NOT NULL},
    nên \texttt{SET NULL} không áp dụng được). Nhóm đã thử đặt \texttt{ON DELETE CASCADE} cho
    \texttt{requests.requested\_by}: SQL Server từ chối với lỗi 1785 (\textit{may cause cycles
    or multiple cascade paths}), vì khi xoá một tài khoản, \texttt{selections} sẽ bị tác động
    theo hai đường: trực tiếp qua \texttt{selections.user\_id} (\texttt{SET NULL}) và gián tiếp
    qua \texttt{requests} (\texttt{CASCADE}).
\end{itemize}

Toàn bộ 24 ràng buộc CHECK (lấy từ \texttt{sys.check\_constraints}):

{\small
\begin{longtable}{>{\raggedright\arraybackslash}p{0.33\textwidth}>{\raggedright\arraybackslash}p{0.6\textwidth}}
\caption{Các ràng buộc CHECK}\label{tab:checks}\\
\toprule \textbf{Tên} & \textbf{Điều kiện} \\ \midrule \endfirsthead
\toprule \textbf{Tên} & \textbf{Điều kiện} \\ \midrule \endhead
\bottomrule \endfoot
\texttt{CK\_asr\_segments\_span} & \texttt{[start\_sec]>=(0) AND [end\_sec]>=[start\_sec]} \\
\texttt{CK\_asr\_segments\_text} & \texttt{len([asr\_text])>(0)} \\
\texttt{CK\_collections\_access} & \texttt{[access\_level]='restricted' OR [access\_level]='open'} \\
\texttt{CK\_eval\_requests\_latency} & \texttt{[latency\_ms]>=(0)} \\
\texttt{CK\_eval\_results\_rank} & \texttt{[rank\_no]>(0)} \\
\texttt{CK\_eval\_runs\_values} & \texttt{[top\_k]>(0) AND [tolerance\_sec]>=(0)} \\
\texttt{CK\_ground\_truth\_span} & \texttt{[start\_sec]>=(0) AND [end\_sec]>=[start\_sec]} \\
\texttt{CK\_keyframes\_values} & \texttt{[pts\_time]>=(0) AND [frame\_idx]>=(0) AND [sample\_n]>(0)} \\
\texttt{CK\_ocr\_texts\_text} & \texttt{len([ocr\_text])>(0)} \\
\texttt{CK\_organizations\_type} & \texttt{[org\_type]='museum' OR [org\_type]='school' OR [org\_type]='press' OR [org\_type]='tv' OR [org\_type]='library'} \\
\texttt{CK\_requests\_origin} & \texttt{[origin]='user' OR [origin]='sample'} \\
\texttt{CK\_requests\_purpose} & \texttt{[purpose]='exhibition' OR [purpose]='education' OR [purpose]='research' OR [purpose]='broadcast'} \\
\texttt{CK\_requests\_status} & \texttt{[status]='rejected' OR [status]='approved' OR [status]='pending'} \\
\texttt{CK\_requests\_text} & \texttt{len([request\_text])>(0)} \\
\texttt{CK\_search\_log\_values} & \texttt{[top\_k]>(0) AND [result\_count]>=(0) AND [latency\_ms]>=(0)} \\
\texttt{CK\_sessions\_expiry} & \texttt{[expires\_at]>[created\_at]} \\
\texttt{CK\_terms\_length} & \texttt{len([term])>=(2)} \\
\texttt{CK\_topics\_parent} & \texttt{[parent\_topic\_id] IS NULL OR [parent\_topic\_id]<>[topic\_id]} \\
\texttt{CK\_usage\_licenses\_period} & \texttt{[valid\_to]>[valid\_from]} \\
\texttt{CK\_users\_type} & \texttt{[user\_type]='other' OR [user\_type]='student' OR [user\_type]='researcher' OR [user\_type]='journalist' OR [user\_type]='staff'} \\
\texttt{CK\_users\_username} & \texttt{len([username])>=(3) AND len([username])<=(32) AND NOT ([username]) collate Latin1\_General\_100\_BIN2 like '\%[\textasciicircum{}a-z0-9\_.-]\%'} \\
\texttt{CK\_videos\_fps} & \texttt{[fps] IS NULL OR [fps]>(0)} \\
\texttt{CK\_videos\_id} & \texttt{[video\_id] like 'L[0-9][0-9][\_]V[0-9][0-9][0-9]'} \\
\texttt{CK\_videos\_title} & \texttt{len([title])>(0)} \\
\end{longtable}
}

\section{Chỉ mục}

Mỗi khoá chính tạo một chỉ mục cụm (clustered) và mỗi ràng buộc UNIQUE tạo một chỉ mục
không cụm. Ngoài ra có 12 chỉ mục \texttt{IX\_} bổ sung, chủ yếu trên các cột khoá ngoại
hay được dùng để nối bảng và lọc. Bảng~\ref{tab:indexes} liệt kê các chỉ mục được những truy
vấn thường gặp sử dụng; hiệu quả được đo ở mục~\ref{sec:tests}.

\begin{table}[H]
\centering
\caption{Các chỉ mục phục vụ truy vấn thường xuyên}
\label{tab:indexes}
\renewcommand{\arraystretch}{1.2}
\begin{tabularx}{\textwidth}{>{\raggedright\arraybackslash}p{6.2cm} >{\raggedright\arraybackslash}X}
\toprule
\textbf{Chỉ mục} & \textbf{Phục vụ} \\
\midrule
\texttt{UQ\_keyframes\_video\_time} (video\_id, pts\_time) & Khung hình gần một thời điểm,
khung hình lân cận (vừa là ràng buộc, vừa là chỉ mục) \\
\texttt{PK\_ocr\_postings} (term\_id, keyframe\_id) & Tra danh sách khung hình của một từ khoá \\
\texttt{UQ\_terms\_term} (term) & Tra mã của từ khoá \\
\texttt{IX\_asr\_segments\_video\_time} (video\_id, start\_sec) & Lời nói quanh một thời điểm \\
\texttt{IX\_requests\_status} (status) & Hàng đợi yêu cầu chờ xử lý \\
\texttt{IX\_requests\_user} (requested\_by) & Bạn đọc xem yêu cầu của mình \\
\texttt{IX\_usage\_licenses\_valid\_to} (valid\_to) & Giấy phép sắp hết hạn \\
\texttt{IX\_videos\_collection}, \texttt{IX\_collections\_organization} & Thống kê theo bộ sưu tập, theo đơn vị \\
\texttt{IX\_topics\_parent}, \texttt{IX\_video\_topics\_topic} & Đi xuống cây chủ đề, tìm video theo chủ đề \\
\texttt{IX\_sessions\_expires} (expires\_at) & Dọn phiên hết hạn \\
\texttt{IX\_audit\_log\_created} (created\_at) & Xem nhật ký mới nhất \\
\bottomrule
\end{tabularx}
\end{table}
